\section{Introduction}
Electroencephalographic event-related potentials (ERPs) are time-locked measurements of neural response dynamics. In biomedical data analysis, they offer millisecond-scale temporal information but are difficult to model because the observed waveform is a noisy, low-dimensional projection of an unobserved biological dynamical system. Inter-subject variability, trial averaging, volume conduction, and measurement noise can cause endpoint classifiers to conflate temporal structure with subject-specific nuisance variation.

Spiking reservoirs provide one way to construct temporally structured representations without replacing signal analysis by opaque endpoint classification. A fixed recurrent reservoir can act as a nonlinear temporal expansion of the input signal, while spike-count summaries provide discrete observables traceable to post-stimulus time windows. The central question is not whether a spiking reservoir universally outperforms classical ERP features, but which temporal information it preserves, transforms, or suppresses under perturbation.

This paper studies a binned spiking reservoir encoding for affective ERP signal analysis. Each ERP observation is mapped through a fixed leaky integrate-and-fire (LIF) reservoir, summarized by six post-onset spike-count bins, and compressed into a reservoir embedding. Fig.~\ref{fig:pipeline} summarizes the pipeline. The contribution is threefold: (i) formulation of affective ERP analysis as reservoir temporal encoding; (ii) subject-level comparison against classical ERP-derived features using balanced accuracy and macro-F1; and (iii) characterization of the encoding under amplitude noise, temporal jitter, and channel dropout.

\begin{figure}[t]
\centering
\resizebox{\linewidth}{!}{% Figure source: reservoir temporal-encoding pipeline for BIBM 2026.
% Compile through main_bibm2026.tex with TikZ loaded.
\begin{tikzpicture}[
    font=\footnotesize,
    >=Latex,
    block/.style={draw, rounded corners, align=center, minimum height=7mm, text width=0.84\linewidth, inner sep=3pt},
    arrow/.style={-{Latex[length=2mm]}, thick}
]
\node[block] (x) {ERP observation\\$X_s^{(c)}\in\mathbb{R}^{M\times T}$};
\node[block, below=4.5mm of x] (lif) {fixed LIF reservoir dynamics\\$v_{t+1}=\beta v_t+Wz_t+W_{\mathrm{in}}x_t-\vartheta r_t$};
\node[block, below=4.5mm of lif] (spikes) {binary spike trajectory\\$z_t=\mathbf{1}[v_t\geq\vartheta]$};
\node[block, below=4.5mm of spikes] (bins) {six-bin spike-count code\\$b_{j,k}=\sum_{t\in\mathcal{B}_k}z_j(t)$};
\node[block, below=4.5mm of bins] (embed) {PCA-compressed reservoir embedding\\$E_s^{(c)}$};
\node[block, below=4.5mm of embed] (eval) {subject-level readout\\clean and perturbed evaluation};

\draw[arrow] (x) -- (lif);
\draw[arrow] (lif) -- (spikes);
\draw[arrow] (spikes) -- (bins);
\draw[arrow] (bins) -- (embed);
\draw[arrow] (embed) -- (eval);
\end{tikzpicture}
}
\caption{Reservoir temporal-encoding pipeline. The object of study is the binned spike-count representation induced by fixed LIF reservoir dynamics, not an end-to-end clinical decision model.}
\label{fig:pipeline}
\end{figure}
